% Section 1 --- Introduction.
\section{Introduction}
\label{sec:intro}

One language model checking another is now the default instrument of oversight.
Judges grade benchmarks, filter pretraining corpora, supply reward signal, and
watch deployed behaviour, and the appeal is obvious: a model is cheap, fast, and
available in quantity where a human annotator is none of those things
\citep{zheng2023judge}. The proposals that go beyond a single judge --- debate
\citep{irving2018debate}, sandwiching \citep{bowman2022measuring},
weak-to-strong training \citep{burns2023weak} --- all inherit the same
arrangement: a model is held to account by other models.

Every version of this rests on an assumption that is rarely stated and almost
never tested. The checker has to fail \emph{differently} from the model it is
checking. If the two share a mistake, no amount of checking surfaces it. And a
checker that reads outputs is blind in two separate ways. It cannot see what the
model keeps to itself, and it cannot see a mistake the two of them make
together.

The first blindness has a known answer. Linear probes read a model's internal
state rather than its output, recovering a direction along which true and false
claims separate \citep{burns2023ccs,marks2024geometry}. The probe literature has
established that these directions exist, that they can be more accurate than the
model's own stated answer \citep{liu2023dissonance}, and that the unsupervised
version of the method is fragile in specific ways
\citep{farquhar2024challenges}. A probe therefore gives a checker access to what
the model represents rather than what it chose to say.

Whether probes answer the second blindness is a different question, and it is the
one we ask. The output-level version has been settled, and the answer is
discouraging: model mistakes are substantially correlated, more so as models get
more capable, and this undermines oversight schemes built on agreement
\citep{goel2025alike,kim2025correlated}. But correlated \emph{outputs} do not
have to mean correlated \emph{representations}. Two models might converge on the
same answer through different internal routes, in which case a reader of
internal state would see the disagreement an output-level judge misses. Nothing
in the existing results rules that out, because nothing has measured it.

\paragraph{False agreement.}
We fit a probe on each of two models, read both on the same claim, and ask how
often they are both wrong in the same direction: both read a claim as true when
the claim is false. We call this \emph{false agreement}. It is the case
disagreement-based oversight cannot see --- the two readouts match, nothing is
flagged, and the answer is wrong.

Its rate is not a free parameter. Restricting attention to the false claims, it
decomposes exactly into a term set by how accurate the two probes are and a term
set by how correlated their mistakes are, and it can never exceed the error rate
of the more accurate probe alone (Section~\ref{sec:theory}). Accuracy enters only
the first term. Two accurate checkers that fail together are therefore worse than
two mediocre ones that fail apart, which is the sense in which the quantity to
select for is not accuracy.

\paragraph{What we do.}
We measure false agreement across $49$ ordered model pairs drawn from ten
open-weight models spanning four families and $1.7$B to $32$B parameters, on five
datasets, with the overseer never larger than the model it oversees. Because a
probe validated on one model is only useful elsewhere if it survives the trip, we
also fit a linear map between each pair's activation spaces and report the cost
of transport explicitly rather than assuming it away. Finally we ask the question
the decomposition raises: if the blind spot is made of shared error, does shared
error leave a signature in the activations that a classifier can learn --- so
that the method could detect its own blind spot?

\paragraph{Findings.}
Three, in descending order of how much we expected them.

First, false agreement is substantial: $4$--$23\%$ of false claims depending on
the dataset, with the two probes' errors positively correlated in nearly every
pair we measure. The blind spot is real at the level of internal belief, not only
at the level of outputs.

Second, and against the obvious prediction, it does \emph{not} shrink when the
two models come from different families. The decomposition explains why: pairs
from different families do have less correlated errors, but they arrive with
higher individual error rates, and the two terms cancel. Choosing a checker from
a different lab is not the lever it appears to be.

Third, the blind spot is largely not visible in the activations that produced it.
A direction fitted on the false agreements of four datasets does transfer to a
fifth, but on the dataset where false agreement is most common it recovers only
five further percentage points of the overseen model's errors beyond what
disagreement already catches. It works well only where false agreement is already
rare. This is the activation-level counterpart of a result proved at the level of
judge scores, that common-mode error is not identifiable from the judges' own
outputs \citep{afrin2026ensembles}.

Taken together these bound how much oversight agreement between two models can
provide --- roughly half of the overseen model's errors in our setting --- and
locate the limit in shared error rather than in the accuracy of either checker.
